\scanpage{50}
\begin{center}4 - 40 -\hfill n\textsuperscript{o} 392\end{center}

\uline{indépendants sur} $k(\underline{p})$ ; \uline{pour qu'ils soient linéairement indépendants sur} $k(\underline{p})$ \uline{pour tout} $\underline{p}\in P$, \uline{il faut et il suffit que} $M/L$ \uline{soit sans torsion}.

Soient $x_{m+1},\ldots,x_n$ des éléments de $M$ qui, avec $x_1,\ldots,x_m$, forment une base de $V$, et soit $N$ le sous-$A$-module libre de $M$ engendré par les $x_i$, $(1\le i\le n)$. Il résulte du n\textsuperscript{o} 3, th. 3 que l'on a $N_{\underline{p}}=M_{\underline{p}}$ pour presque tout $\underline{p}\in P$ ; comme $x_1(\underline{p}),\ldots,x_n(\underline{p})$ forment une base de $N(\underline{p})$ sur $k(\underline{p})$, cela établit la première assertion.

Si $M/L$ est sans torsion, il en est de même \unsure{dc|de} $(M/L)_{\underline{p}}=M_{\underline{p}}/L_{\underline{p}}$ pour tout $\underline{p}\in P$, (n\textsuperscript{o} 1, \uline{Exemple} 6), et comme $A_{\underline{p}}$ est principal, $M_{\underline{p}}/L_{\underline{p}}$ est libre. Par suite, $M_{\underline{p}}$ est somme directe de $L_{\underline{p}}$ et d'un $A_{\underline{p}}$-module libre $E$ de rang $n-m$ ; donc $M(\underline{p})$ est somme directe de $L(\underline{p})$ et du $k(\underline{p})$-espace vectoriel $E/\underline{p}E$, de dimension $n-m$ ; par suite $L(\underline{p})$ est de dimension $m$, et comme il est engendré par les $x_i(\underline{p})$, $(1\le i\le m)$, ces derniers sont linéairement indépendants.

Inversement, supposons que les $x_i(\underline{p})$, $(1\le i\le m)$, soient linéairement indépendants sur $k(\underline{p})$ pour tout $\underline{p}\in P$. Alors $L_{\underline{p}}$ est facteur direct de $M_{\underline{p}}$ pour tout $\underline{p}$, (chap. II, § 3, n\textsuperscript{o} 2, cor. 1 et 2 de la prop. 5), et par suite $M_{\underline{p}}/L_{\underline{p}}=(M/L)_{\underline{p}}$ est sans torsion pour tout $\underline{p}\in P$. On en conclut que $P\cap\operatorname{Ass}(M/L)=\varnothing$ en vertu du chap. IV, § 1, n\textsuperscript{o} 2, cor. de la prop. 5. Mais comme $L$ est réflexif, il résulte du n\textsuperscript{o} 2, prop. 6, (i), que le seul idéal premier qui puisse appartenir à $\operatorname{Ass}(M/L)$ est l'idéal $(0)$ ; donc $M/L$ est sans torsion.

LEMME 8.- \uline{Supposons que le rang} $n$\unsure{-}\uline{de} $M$ \uline{soit} $\ge2$ ; \uline{alors il existe un élément} $x\ne0$ \uline{de} $M$ \uline{tel que} $M/Ax$ \uline{soit sans torsion}.

Soit $y\ne0$ un élément de $M$. En vertu du lemme 7, l'ensemble $Y$ des $\underline{p}\in P$ tels que $y(\underline{p})=0$ est fini. Si $Y=\varnothing$, il résulte du lemme 7 appliqué à la suite $(x_i)$ formée du seul élément $y$, que $M/Ay$ est sans torsion. Supposons donc $Y\ne\varnothing$, et posons $S=\bigcap_{\underline{p}\in Y}(A-\underline{p})$ ; on sait, (n\textsuperscript{o} 4, lemme 2), que $S^{-1}A$ est un anneau

\margin{49}
